%!TEX root = ../template.tex
%%%%%%%%%%%%%%%%%%%%%%%%%%%%%%%%%%%%%%%%%%%%%%%%%%%%%%%%%%%%%%%%%%%%
%% appendix-cascade.tex
%% Resolution ladder sequence diagrams (ported from ISEL-Thesis-2526/thesis/appendix.tex)
%%%%%%%%%%%%%%%%%%%%%%%%%%%%%%%%%%%%%%%%%%%%%%%%%%%%%%%%%%%%%%%%%%%%
\chapter{Complete Diagrams}
\label{app:cascade-sequences}

This appendix gives the complete diagrams that Chapter~\ref{cha:impl} shows in part: the class
model of the integration layer (Figure~\ref{fig:omniflow-integration-classes}), presented in four
parts in Section~\ref{sec:function-registry}; the
AWS deployment sequence (Figure~\ref{fig:aws-deploy-sequence}), whose Level~3 branch is
Figure~\ref{fig:aws-deploy-level3}, and one diagram per level of the resolution ladder
(Section~\ref{sec:resolution-cascade}), with the GCP path as panel (a) and the AWS path as panel (b).

% The diagrams are placed in sequence rather than as floats, each at the largest size the page
% allows; the two panels of Figures A.2 and A.3 are too tall to share a page, so (b) continues the figure.
\newcommand{\appdiagram}[2][0.85]{%
  \includegraphics[width=\textwidth,height=#1\textheight,keepaspectratio]{images/implementation/#2.png}}

\noindent\begin{minipage}{\textwidth}
  \centering
  \appdiagram[0.75]{omniflow-integration-classes}
  \captionof{figure}{Class model of the integration layer added to OmniFlow.}
  \label{fig:omniflow-integration-classes}
\end{minipage}

\bigskip
\noindent\begin{minipage}{\textwidth}
  \centering
  \appdiagram[0.8]{aws-deploy-sequence}
  \par (a) Overview
  \captionof{figure}{Deployment sequence on AWS, from the registry bootstrap to
  \texttt{Create\allowbreak State\allowbreak Machine}, with one branch per level of the resolution
  ladder and a fourth for the calls that fail fast (a); the Level~3 branch is expanded in (b), on
  the next page.}
  \label{fig:aws-deploy-sequence}
\end{minipage}

\bigskip
\noindent\begin{minipage}{\textwidth}
  \centering
  \appdiagram[0.8]{aws-deploy-sequence-level3-full}
  \par (b) Level~3 expanded
  \captionsetup{type=figure}\ContinuedFloat
  \caption[]{Deployment sequence on AWS, continued: the Level~3 branch, from
  \texttt{deploy\allowbreak OrUpdate} to the \texttt{Function\allowbreak Invocation\allowbreak
  Metadata} it returns.}
\end{minipage}

\bigskip
\noindent\begin{minipage}{\textwidth}
  \centering
  \appdiagram[0.84]{resolution-cascade-level1-gcp}
  \par (a) GCP\par
  \captionof{figure}{Level~1 (Registry): a registry hit is validated against the live Cloud Run
  service on GCP (a) and against \texttt{GetFunction} on AWS (b, on the next page). On a
  \texttt{NotFound} result, both paths attempt rediscovery (across every region, or in the one
  region a \texttt{"region/functionRef"} reference names) before a stale entry is removed.}
  \label{fig:resolution-cascade-level1}
\end{minipage}

\bigskip
\noindent\begin{minipage}{\textwidth}
  \centering
  \appdiagram[0.84]{resolution-cascade-level1-aws}
  \par (b) AWS
  \captionsetup{type=figure}\ContinuedFloat
  \caption[]{Level~1 (Registry), continued: the AWS path.}
\end{minipage}

\bigskip
\noindent\begin{minipage}{\textwidth}
  \centering
  \appdiagram[0.4]{resolution-cascade-level2-gcp}
  \par (a) GCP\par
  \medskip
  \appdiagram[0.4]{resolution-cascade-level2-aws}
  \par (b) AWS
  \captionof{figure}{Level~2 (Provider discovery): the target Cloud Run service is searched for
  across the project's regions on GCP (a) and the target Lambda across every region enabled for the
  account on AWS (b). On both providers, the fully-qualified \texttt{"region/functionRef"} form is
  resolved with a single lookup in that region.}
  \label{fig:resolution-cascade-level2}
\end{minipage}

\bigskip
\noindent\begin{minipage}{\textwidth}
  \centering
  \appdiagram[0.4]{resolution-cascade-level3-gcp}
  \par (a) GCP\par
  \medskip
  \appdiagram[0.4]{resolution-cascade-level3-aws}
  \par (b) AWS
  \captionof{figure}{Level~3 (QuickFaaS deployment): \texttt{QuickFaasDeployer} invokes QuickFaaS,
  polls the Cloud Functions~v2 API until the function is active and grants the workflow's service
  account \texttt{roles/\allowbreak run.invoker} on its Cloud Run service on GCP (a).
  \texttt{AwsLambdaDeployer} invokes QuickFaaS and polls \texttt{GetFunction} until the Lambda is
  active on AWS (b). The AWS branch is expanded in Figure~\ref{fig:aws-deploy-sequence}.}
  \label{fig:resolution-cascade-level3}
\end{minipage}
